% TODO proof-reading
\subsubsection{事前初期化された文字列配列}
\label{Java_2D_array_month}

\begin{lstlisting}[style=customjava]
class Month
{
	public static String[] months = 
	{
		"January", 
		"February", 
		"March", 
		"April",
		"May",
		"June",
		"July",
		"August",
		"September",
		"October",
		"November",
		"December"
	};

	public String get_month (int i)
	{
		return months[i];
	};
} 
\end{lstlisting}

\TT{get\_month()}関数は単純です：

\begin{lstlisting}
  public java.lang.String get_month(int);
    flags: ACC_PUBLIC
    Code:
      stack=2, locals=2, args_size=2
         0: getstatic     #2         // Field months:[Ljava/lang/String;
         3: iload_1       
         4: aaload        
         5: areturn       
\end{lstlisting}

\TT{aaload}は\IT{参照}の配列で動作します。

Java文字列はオブジェクトなので、それらを操作するには\IT{a}命令が使用されます。

\TT{areturn}は\TT{String}オブジェクトへの\IT{参照}を返します。


\TT{months[]}配列はどのように初期化されるのでしょうか？


\begin{lstlisting}
  static {};
    flags: ACC_STATIC
    Code:
      stack=4, locals=0, args_size=0
         0: bipush        12
         2: anewarray     #3         // class java/lang/String
         5: dup           
         6: iconst_0      
         7: ldc           #4         // String January
         9: aastore       
        10: dup           
        11: iconst_1      
        12: ldc           #5         // String February
        14: aastore       
        15: dup           
        16: iconst_2      
        17: ldc           #6         // String March
        19: aastore       
        20: dup           
        21: iconst_3      
        22: ldc           #7         // String April
        24: aastore       
        25: dup           
        26: iconst_4      
        27: ldc           #8         // String May
        29: aastore       
        30: dup           
        31: iconst_5      
        32: ldc           #9         // String June
        34: aastore       
        35: dup           
        36: bipush        6
        38: ldc           #10        // String July
        40: aastore       
        41: dup           
        42: bipush        7
        44: ldc           #11        // String August
        46: aastore       
        47: dup           
        48: bipush        8
        50: ldc           #12        // String September
        52: aastore       
        53: dup           
        54: bipush        9
        56: ldc           #13        // String October
        58: aastore       
        59: dup           
        60: bipush        10
        62: ldc           #14        // String November
        64: aastore       
        65: dup           
        66: bipush        11
        68: ldc           #15        // String December
        70: aastore       
        71: putstatic     #2         // Field months:[Ljava/lang/String;
        74: return        
\end{lstlisting}

\TT{anewarray}は\IT{参照}の新しい配列を作成します（したがって\IT{a}接頭辞）。

オブジェクトの型は\TT{anewarray}のオペランドで定義され、それは\\
\q{java/lang/String}文字列です。

\TT{anewarray}の前の\TT{bipush 12}は配列のサイズを設定します。

ここで新しい命令が見えます：\TT{dup}。


\myindex{Forth}
これはスタックコンピュータ（Forthプログラミング言語を含む）の標準命令で、\ac{TOS}の値を単純に複製します。

\myindex{x86!\Instructions!FDUP}
ちなみに、FPU 80x87もスタックコンピュータで、類似の命令を持っています -- \INS{FDUP}。


これは配列への\IT{参照}を複製するために使用されます。なぜなら\TT{aastore}命令がスタックから配列への\IT{参照}をポップしますが、次の\TT{aastore}が再びそれを必要とするからです。

Javaコンパイラは、各配列ストア操作の前に\TT{getstatic}命令を生成する（つまり11回）よりも、\TT{dup}を生成する方が良いと判断しました。


\TT{aastore}は\ac{TOS}から取得されたインデックスで配列に\IT{参照}（文字列への）を格納します。


最後に、\TT{putstatic}は新しく作成された配列への\IT{参照}をオブジェクトの2番目のフィールド、つまり\IT{months}フィールドに格納します。